\section{Руководство пользователя}
\label{sec:manual}

Для работы с системой не нужно устанавливать отдельное приложение: достаточно клиента \textit{Telegram} на смартфоне, компьютере или в браузере. Пользователь находит бота по имени и отправляет команду /\textit{start}. Отдельной регистрации нет. При первом сообщении система сама создаёт учётную запись по идентификатору \textit{Telegram}, а язык интерфейса выбирает по языку приложения \textit{Telegram}: русский или английский.

В ответ на команду /\textit{start} бот присылает приветствие с кратким описанием работы и кнопки главного меню: «Мои товары», «Добавить товар», «Настройки» и «Помощь». Те же четыре кнопки закрепляются под полем ввода сообщения, так что главное меню доступно на любом экране. Главное меню показано на рисунке~\ref{fig:manual:start}.

\begin{figure}[H]
	\centering
	\includegraphics[scale=0.6]{kufar/screens/start.png}
	\caption{Приветствие и главное меню бота}
	\label{fig:manual:start}
\end{figure}

Кроме кнопок, бот понимает команды: /\textit{start}~-- главное меню, /\textit{add}~-- добавить товар, /\textit{my}~-- мои товары, /\textit{report}~-- отчёт в \textit{Excel}, /\textit{settings}~-- настройки, /\textit{language}~-- язык интерфейса, /\textit{profile}~-- профиль, /\textit{help}~-- справка. Список команд продублирован в меню \textit{Telegram} слева от поля ввода.

\subsection{Роль администратора}
\label{sec:manual:admin}

Администратор~-- пользователь, идентификатор \textit{Telegram} которого указан в конфигурации системы. Отдельного входа или пароля у администратора нет: бот распознаёт его по учётной записи \textit{Telegram}. Чтобы назначить администратора, владелец сервера добавляет его идентификатор в список администраторов в файле окружения и перезапускает сервисы. Узнать свой идентификатор можно в разделе профиля бота.

Администратор пользуется всеми функциями пользователя, описанными в подразделе~\ref{sec:manual:user}. Возможности администратора отличаются в трёх местах.

\subsubsection{Повышенный лимит отслеживания}

Администратор может отслеживать до 10 товаров вместо 3. В профиле (команда /\textit{profile}) в строке «Статус» указано «администратор», а в строке «Лимит отслеживания» видно, сколько товаров занято у него самого и во всей системе. Общий предел в 25 товаров действует и для администратора и защищает площадку от перегрузки запросами. Профиль администратора показан на рисунке~\ref{fig:manual:admin_profile}.

\begin{figure}[H]
	\centering
	\includegraphics[scale=0.6]{kufar/screens/admin_profile.png}
	\caption{Профиль администратора}
	\label{fig:manual:admin_profile}
\end{figure}

\subsubsection{Отчёт о состоянии системы}

Под профилем расположена кнопка «Отчёт в \textit{Excel}». В меню отчётов администратор, помимо кнопок периода, видит кнопку «Состояние системы (30 дн.)». По ней бот присылает книгу \textit{Excel}, в которой собраны показатели за последние 30 дней:
\begin{itemize}
	\item число пользователей всего и активных;
	\item число отслеживаемых товаров и из них активных;
	\item число совпадений, найденных за период, и доставленных уведомлений;
	\item число совпадений, ожидающих доставки, и не доставленных после всех попыток;
	\item распределение отслеживаемых товаров по категориям.
\end{itemize}

По этому отчёту администратор видит, справляется ли система с нагрузкой. Растущее число ожидающих доставки совпадений означает, что уведомления не успевают уходить, а ненулевое число недоставленных говорит о том, что часть пользователей не получает сообщения. Пример отчёта приведён на рисунке~\ref{fig:manual:system_report}.

\begin{figure}[H]
	\centering
	\includegraphics[scale=0.6]{kufar/screens/system_report.png}
	\caption{Отчёт о состоянии системы}
	\label{fig:manual:system_report}
\end{figure}

\subsubsection{Служебные оповещения}

Администратор получает от бота сообщения о событиях, которые требуют внимания человека:
\begin{itemize}
	\item площадка ответила отказом (код 403 или 429)~-- в сообщении указано, на сколько циклов приостановлена проверка, и что делать, если отказ повторяется: снизить общий лимит товаров или увеличить интервал опроса;
	\item контейнер системы запущен, остановлен штатно или упал с ненулевым кодом выхода;
	\item контейнер перестал проходить проверку работоспособности.
\end{itemize}

Одно и то же оповещение не повторяется подряд. Сообщение о блокировке площадкой приходит не чаще раза в 30 минут, о событиях контейнера не чаще раза в 5 минут. Если контейнер перезапускается чаще, в оповещении отдельно указано, что это похоже на цикл перезапусков. Пример оповещения показан на рисунке~\ref{fig:manual:admin_alert}.

\begin{figure}[H]
	\centering
	\includegraphics[scale=0.6]{kufar/screens/admin_alert.png}
	\caption{Служебное оповещение администратору}
	\label{fig:manual:admin_alert}
\end{figure}

\subsubsection{Обслуживание сервера}

Часть задач администратор выполняет не в \textit{Telegram}, а на сервере. Все сервисы запускаются и обновляются одной командой \textit{Docker Compose}, журналы работы доступны там же. Резервные копии базы данных создаются автоматически раз в сутки в отдельном каталоге сервера. Пригодность последней копии проверяется сценарием восстановления, который загружает её во временную базу и выводит число строк в основных таблицах. Письма системы при стандартной настройке попадают в почтовый сервер-ловушку, который доступен в браузере на сервере; для доставки настоящим адресатам в файле окружения указываются параметры внешнего почтового сервера.

\subsection{Роль пользователя}
\label{sec:manual:user}

\subsubsection{Добавление отслеживаемого товара}

Добавление начинается с кнопки «Добавить товар» или команды /\textit{add}. Бот предлагает выбрать категорию из 13 вариантов. Первыми в списке стоят «Велосипед», «Ноутбук» и «Квартира (аренда)»: в них отбор ведётся по фильтрам самого \textit{Kufar.by} и получается заметно точнее. Над списком категорий находится кнопка «Вставить ссылку с \textit{Kufar}»; работа с ней описана далее. Экран выбора категории показан на рисунке~\ref{fig:manual:category}.

\begin{figure}[H]
	\centering
	\includegraphics[scale=0.6]{kufar/screens/category.png}
	\caption{Выбор категории товара}
	\label{fig:manual:category}
\end{figure}

На следующем шаге вводится название товара: бренд и модель, например «\textit{Trek Marlin}» или «\textit{Pixel} 7\textit{a}». Лишних слов писать не стоит: бот требует, чтобы в объявлении встретилось каждое слово названия. В категориях с фильтрами площадки название можно пропустить кнопкой «Пропустить», тогда отбор выполнят фильтры.

Для велосипедов, ноутбуков и квартир бот открывает меню фильтров. В нём перечислены характеристики категории: для велосипеда это производитель, класс, размер рамы, диаметр колёс, материал рамы, для кого, состояние, регион и город. Пользователь открывает характеристику и отмечает одно или несколько значений; объявление подойдёт, если в нём есть любое из отмеченных. Длинные списки, например бренды или процессоры, разбиты на страницы по восемь значений. Кнопка «Сбросить» снимает выбор, «Назад» возвращает к списку характеристик, «Готово» завершает настройку. Сначала лучше выбрать регион: после этого список городов и районов сократится до выбранного региона. Меню фильтров показано на рисунке~\ref{fig:manual:filters}.

\begin{figure}[H]
	\centering
	\includegraphics[scale=0.6]{kufar/screens/filters.png}
	\caption{Выбор значений фильтра}
	\label{fig:manual:filters}
\end{figure}

У квартир есть числовые фильтры: общая площадь, этаж и год постройки. Их значения вводятся сообщением в виде диапазона («45-70»), нижней границы («от 2015») или одного числа. Символ «-» сбрасывает фильтр.

Затем бот просит указать цену в белорусских рублях. Это минимальная цена, по которой товар сейчас продаётся на рынке: бот будет присылать объявления дешевле неё. На следующем шаге бот показывает стоп-слова, по которым объявления отсеиваются автоматически (например, «чехол» для электроники), и предлагает добавить свои слова-исключения через запятую, например «на запчасти, разбит». Шаг можно пропустить.

В конце бот показывает сводку: название, категорию, фильтры, цену и исключения. После подтверждения товар сохраняется, и бот сразу ищет на площадке объявления дешевле указанной цены. До пяти самых дешёвых из них приходят отдельным сообщением. Если таких объявлений нет, бот сообщает об этом и продолжает следить за площадкой. Результат мгновенного поиска показан на рисунке~\ref{fig:manual:instant}.

\begin{figure}[H]
	\centering
	\includegraphics[scale=0.6]{kufar/screens/instant.png}
	\caption{Сохранение товара и мгновенный поиск}
	\label{fig:manual:instant}
\end{figure}

Добавить товар не получится, если занят личный лимит (3 товара) или общий лимит системы (25 товаров). Бот объясняет, какой из них достигнут. В первом случае нужно удалить ненужный товар в разделе «Мои товары», во втором остаётся подождать, пока место освободится.

На любом шаге диалог можно прервать кнопкой «Отмена», введённые данные при этом не сохраняются.

\subsubsection{Импорт фильтров из ссылки}

Если поиск уже настроен на сайте \textit{Kufar.by}, его не нужно повторять в боте. Пользователь настраивает фильтры на сайте или открывает свой сохранённый поиск, копирует ссылку из адресной строки и отправляет её боту после нажатия «Вставить ссылку с \textit{Kufar}». Бот переносит категорию, фильтры и цену и перечисляет неподдерживаемые параметры, например фильтр по станции метро. Цену из ссылки можно оставить одной кнопкой или ввести другую. Ссылку на страницу сохранённого поиска в личном кабинете бот принять не может, потому что её содержимое видно только владельцу, и подсказывает, какую ссылку прислать вместо неё.

\subsubsection{Получение уведомлений}

Когда на площадке появляется подходящее объявление, бот присылает сообщение. В заголовке указан тип события: «Новое предложение» или «Цена упала». В сообщении приведены цена (при снижении~-- старая зачёркнутая и новая), на сколько процентов цена ниже порога пользователя, состояние товара, местоположение, время публикации и ссылка «Открыть объявление». Первая фотография объявления идёт вместе с текстом, ещё до четырёх приходят следом альбомом. Пример уведомления показан на рисунке~\ref{fig:manual:notification}.

\begin{figure}[H]
	\centering
	\includegraphics[scale=0.6]{kufar/screens/notification.png}
	\caption{Уведомление о выгодном предложении}
	\label{fig:manual:notification}
\end{figure}

Под уведомлением расположены кнопки «Интересно», «Не то» и «Куплю». Оценка сохраняется и попадает в отчёт, а число покупок по рекомендации бота видно в профиле. Одно и то же объявление повторно приходит только в том случае, если его цена стала ниже, чем в прошлом уведомлении. Чтобы уведомлений не было слишком много, по одному товару приходит не более 5 сообщений в час и не более 30 в сутки на пользователя; при большом числе находок первыми приходят самые дешёвые.

\subsubsection{Управление списком товаров}

Кнопка «Мои товары» или команда /\textit{my} открывает список отслеживаемых товаров. Над списком указано, сколько товаров отслеживается из личного лимита и сколько занято во всей системе. Под списком находятся кнопки «Экспорт \textit{CSV}» и «Импорт \textit{CSV}». Нажатие на товар открывает его карточку с категорией, ценой, фильтрами, статусом, словами-исключениями и автоматическими стоп-словами. Карточка товара показана на рисунке~\ref{fig:manual:product}.

\begin{figure}[H]
	\centering
	\includegraphics[scale=0.6]{kufar/screens/product.png}
	\caption{Карточка отслеживаемого товара}
	\label{fig:manual:product}
\end{figure}

В карточке доступны действия:
\begin{itemize}
	\item «Редактировать»~-- изменить название, цену, слова-исключения, а в категориях с фильтрами площадки и сами фильтры;
	\item «Пауза/Возобновить»~-- временно остановить отслеживание без удаления товара;
	\item «Удалить»~-- прекратить отслеживание и удалить товар;
	\item «Назад»~-- вернуться к списку.
\end{itemize}

Кнопка «Экспорт \textit{CSV}» присылает файл со списком товаров, который сразу открывается в \textit{Excel} по столбцам. Файл можно отредактировать и загрузить обратно кнопкой «Импорт \textit{CSV}», например чтобы перенести список на другую учётную запись. При импорте бот проверяет каждую строку так же, как при ручном вводе, и сообщает, сколько товаров добавлено, какие строки пропущены и почему.

\subsubsection{Настройки и электронная почта}

Кнопка «Настройки» или команда /\textit{settings} показывает, включены ли уведомления, выбранный язык и состояние электронной почты. Здесь можно отключить и снова включить уведомления, сменить язык интерфейса и перейти в раздел «Почта».

В разделе «Почта» пользователь нажимает «Привязать адрес» и вводит адрес электронной почты. На него приходит шестизначный код, который действует 15 минут; его нужно ввести в бота. Если письмо не пришло, код можно запросить снова, но не чаще раза в минуту. После подтверждения копии уведомлений раз в 5 минут приходят на почту сводным письмом, а отчёты можно получать вложением. Копии отключаются кнопкой в том же разделе, там же адрес можно сменить или отвязать.

\subsubsection{Профиль и отчёты}

Команда /\textit{profile} показывает идентификатор, имя, статус, число отслеживаемых и активных товаров, найденных предложений, покупок по рекомендации и занятость лимита. Кнопка «Отчёт в \textit{Excel}» под профилем или команда /\textit{report} предлагает выбрать период: 7, 30 или 90 дней. При подтверждённой почте появляются такие же кнопки для отправки отчёта письмом.

Отчёт состоит из двух листов. На листе «Находки» по одной строке на каждое найденное предложение: дата, товар, объявление, цена, порог, выгода в процентах, тип события, статус доставки, оценка и ссылка. Цены записаны числами, поэтому таблицу можно сортировать и фильтровать средствами \textit{Excel}. На листе «Итоги» указаны период, число находок и доставленных уведомлений, средняя выгода и лучшая находка. Пример отчёта приведён на рисунке~\ref{fig:manual:report}.

\begin{figure}[H]
	\centering
	\includegraphics[scale=0.6]{kufar/screens/report.png}
	\caption{Отчёт о найденных предложениях}
	\label{fig:manual:report}
\end{figure}

\subsubsection{Действия при ошибках}

Бот отвечает на каждое действие. Если сервис временно недоступен, пользователь получает сообщение с просьбой повторить попытку через минуту. Если бот не понял сообщение, например диалог добавления прервался после перезапуска, пользователю предлагается начать заново кнопкой «Добавить». Команда /\textit{help} в любой момент показывает список команд и краткое описание работы бота.
